\section{The split $E_8$ rational surface}

\subsection{Why torsion-friendly models are too small}

A rational elliptic surface with rational two-torsion can be written
\[
y^2=x^3+A(t)x^2+B(t)x,
\qquad \deg A\le2,\quad \deg B\le4.
\]
Its discriminant is
\[
\Delta=16B^2(A^2-4B).
\]
The zeros of $B$ force reducible fibres of total root rank at least four, so
Shioda--Tate bounds the generic Mordell--Weil rank by four.  A smooth
three-character $V_4$ packet built from such bases has total capacity at most
twenty-two.

Rational three-torsion in Tate normal form forces still more reducible-fibre
cost; the analogous packet capacity is at most fourteen.  The easiest descent
models are therefore incapable of carrying rank thirty.  The high-capacity
base should have trivial small rational torsion and twelve irreducible
$I_1$ fibres.

\subsection{A certified split $E_8$ surface}

The repository contains an exact certificate for the rational elliptic
surface
\[
\mathcal E:\quad y^2=x^3+A(t)x+B(t),
\]
with
\[
\begin{aligned}
A(t)={}&-27t^4+216t^3+6t^2-432t-267,\\
B(t)={}&54t^6+108t^5-630t^4-880t^3+1458t^2+1278t+1242.
\end{aligned}
\]
Its discriminant has twelve simple roots, so every singular fibre is of type
$I_1$.  The surface is rational and has no reducible-fibre root lattice.
Eight explicit rational sections are stored in the machine-readable
certificate.  Their Shioda height matrix is even, positive definite, has
rank eight, and determinant one.  Therefore the generic Mordell--Weil lattice
is exactly
\[
\boxed{E_8.}
\]

This is the preferred base for packet engineering: it realizes the maximum
rank of a rational elliptic surface without spending Picard rank on reducible
fibres.

\subsection{Why higher-$\chi$ neighbour fibrations do not replace K3 methods}

Let $S\to\mathbf P^1$ be a relatively minimal elliptic surface with section
and $d=\chi(\mathcal O_S)$.  The canonical bundle formula gives
\[
K_S\sim(d-2)F.
\]
If $d\ne2$, the canonical class determines the primitive fibre ray.  Any
other genus-one fibre class $F'$ satisfies $K_S\cdot F'=0$, hence
$F\cdot F'=0$; the Hodge index theorem forces $F'$ to be proportional to
$F$.  Primitivity gives $F'=F$.

Thus
\[
\boxed{d\ne2\Longrightarrow\text{the genus-one fibration is canonically rigid}.}
\]
K3 surfaces, where $d=2$ and $K_S=0$, are the exceptional case in which
neighbour fibrations can genuinely move lattice directions between fibre
components and sections.  This explains why changing the K3 fibration broke
the rank-$28$ plateau, while the same manoeuvre cannot simply be transplanted
to an arithmetic-genus-four surface.
